\begin{claim}{D.64}{Proved}{Standard negative torus centralizers and their actual kernels}
Let $b_0,b_1$ be the two standard negative outer
involutions on $F_0=\widehat F(X,Y)$:
\[
 \begin{array}{c|cc}
       &X&Y\\ \hline
 \alpha_0&YXY^{-1}&Y^{-1}\\
 \alpha_1&Y&X .
 \end{array}                                                  \tag{D.64a}
\]
Both have the named cusp $c_0=[X,Y]$ and
$\alpha_j(c_0)=c_0^{-1}$.
For $b_0$, put $\zeta_0(X)=X^{-1}$,
$\zeta_0(Y)=Y$.
Its full peripheral outer centralizer and normalizer
are exactly
\[
 C_\Omega(b_0)=N_\Omega(\langle b_0\rangle)
      =\langle b_0,[\zeta_0]\rangle\cong C_2^2.               \tag{D.64b}
\]
Thus every fibre of the actual orbit map
$\theta\mapsto\theta b_0\theta^{-1}$ is precisely
a right coset of this specified four-element group.
Conjugation by $\theta$ gives the exact corresponding
centralizer on that orbit.

For both $j=0,1$, the IA kernel $I_{b_j}$ of
D.63 is trivial. Consequently the full centralizer
of $b_1$ has the faithful character embedding
\[
 C_\Omega(b_1)\ \lhook\joinrel\longrightarrow\
          \{1,-1\}\times\widehat\Z^\times,\qquad
                         \beta\longmapsto(\sigma,\lambda).
                                                               \tag{D.64c}
\]
In particular it is abelian.
For the actual automorphism
$\zeta_1(X)=X^{-1}$, $\zeta_1(Y)=Y^{-1}$,
\[
 \begin{split}
 C_\Omega(b_1)\cap\ker\lambda&=\langle[\zeta_1]\rangle
                                                  \cong C_2,\\
 \{\beta\in C_\Omega(b_1):\lambda(\beta)\in\{1,-1\}\}
                   &=\langle b_1,[\zeta_1]\rangle\cong C_2^2.
 \end{split}                                                  \tag{D.64d}
\]
This is also the entire torsion subgroup of this
abelian centralizer. The remaining multiplier
calculation for (D.64c) is supplied by D.65 below.
\end{claim}

\begin{proof}
\emph{The exact standard extensions.}
D.61 gives the mutually inverse full presentation maps
\[
 \begin{aligned}
 E_0&=(\widehat\Z\langle X\rangle\times C_2\langle w\rangle)
                                      \amalg C_2\langle v\rangle,\\
 Y&=vw,\qquad w=Y^{-1}v,\\
 E_1&=C_2\langle v\rangle\amalg\widehat\Z\langle z\rangle,\\
 X&=zv,\qquad Y=vz,\qquad z=Xv .
 \end{aligned}                                                \tag{D.64e}
\]
The parity homomorphisms have
$\epsilon(X)=0,\epsilon(v)=\epsilon(w)=1$ in the
first row and $\epsilon(v)=\epsilon(z)=1$ in the
second. Their kernels are the actual $F_0$ in
(D.64a), and conjugation by $v$ gives $\alpha_j$.
The full free-factor intersection theorem used
in D.35 gives
\[
 C_{E_0}(v)=\langle v\rangle,\qquad
 C_{E_0}(w)=\langle X\rangle\times\langle w\rangle,\qquad
 C_{E_1}(v)=\langle v\rangle.
                                                               \tag{D.64f}
\]
For example a centralizer element outside the
specified factor would put the nontrivial involution
in the intersection of that factor with its conjugate.
The intersection theorem excludes that possibility;
the indicated factor is already abelian where needed
\cite[Theorem 5 of the inspected author version]{A-HerfortZalesskii2008}.
In particular $C_{F_0}(v)=1$ in both rows.

The actual fusion transporters are
\[
          g_0=X,\qquad g_1=z^2=XY,\qquad
          g_jvg_j^{-1}=c_0v.
                                                               \tag{D.64g}
\]
The identities follow by substitution in (D.64e).
Both have primitive homology, and
$c_0=[g_j,v]$ in their respective extensions.

\emph{Vanishing of the two actual IA kernels.}
Use D.63f to represent an IA centralizer element
by its unique actual automorphism fixing $c_0$
and commuting with $\alpha_j$.
Its extension fixes $v$.
The exact multiplier-one transporter formula in
D.62, together with $C_{F_0}(v)=1$, makes it fix
$g_j$ as well.

For $E_0$, D.62 says that $w$ is the unique
outside involution centralizing $g_0=X$.
The extension must fix it.
It fixes all three generators $X,w,v$ of $E_0$,
and is therefore the identity.
For $E_1$, D.62 says that $z$ is the unique
outside square root of $g_1=z^2$.
The extension must fix it, and hence fixes the
two generators $z,v$.
This proves $I_{b_0}=I_{b_1}=1$ using actual
extension generators and their uniqueness, rather
than generation inferred from homology.
The faithful map (D.64c) and abelianness follow
from D.63e.

\emph{All multipliers in the two-class standard model.}
Let $\beta\in C_\Omega(b_0)$.
Normalize its actual extension by D.63:
\[
        \Psi(c_0)=c_0^\mu,\qquad
        \Psi(v)=c_0^iv,\qquad i\in\{0,1\}.
\]
There is no fixed-subgroup ambiguity in its
transporter image, since the two groups $C_i$
in D.62 are trivial by (D.64f).
Thus, with $a=\Psi(X)$,
\[
       a=
       \begin{cases}
         c_0^{(\mu-1)/2}X,&i=0,\\
         c_0^{(\mu+1)/2}X^{-1},&i=1.
       \end{cases}                                          \tag{D.64h}
\]
The involution $t=\Psi(w)$ belongs to the other
full $F_0$-section class: $\Psi$ permutes those
two classes, and already preserves the cusp
class by its normalized image of $v$.
Therefore $t=hwh^{-1}$ for some actual $h\in F_0$.
Since $a$ centralizes $t$, (D.64f) gives
\[
 a\in h\bigl(C_{E_0}(w)\cap F_0\bigr)h^{-1}
       =h\langle X\rangle h^{-1}.
 \]
Write $a=hX^u h^{-1}$.
Its primitive homology in D.62 is $(-1)^i[X]$,
so $u=(-1)^i$ in the full $\widehat\Z$.
This is the exact conjugacy
\[
                       a\sim_{F_0}X^{(-1)^i}.                \tag{D.64i}
\]

For every integer $n\geq2$, use the explicit
finite dihedral quotient
\[
 \begin{split}
 q_n:E_0&\longrightarrow
       D_{2n}=\langle r,s\mid r^n=s^2=1,\ srs=r^{-1}\rangle,\\
       X&\longmapsto r,\qquad v\longmapsto s,\qquad
       w\longmapsto1 .
 \end{split}                                                  \tag{D.64j}
\]
It respects the defining relations in (D.64e)
and is surjective.
The exact peripheral image is $q_n(c_0)=r^2$.
Both lines of (D.64h) consequently give
$q_n(a)=r^\mu$.
By (D.64i), this rotation is conjugate to
$r^{(-1)^i}$.
The conjugacy class of either $r$ or $r^{-1}$
in this finite group is $\{r,r^{-1}\}$.
Hence
\[
                  \mu\bmod n\in\{1,-1\}
                       \quad\text{for every }n\geq2.
                                                               \tag{D.64k}
\]
These simultaneous congruences force one global
sign $\mu=1$ or $\mu=-1$ in $\widehat\Z$.
For a direct compactness proof, the sets
$\{\varepsilon\in\{1,-1\}:\mu\equiv\varepsilon
\pmod n\}$ are nonempty and have the finite
intersection property, since the condition for
a least common multiple implies every specified
finite set of conditions.
Their intersection supplies the one sign at
all moduli.

Together with $I_{b_0}=1$, D.63 now bounds its
full centralizer by four elements.
All four occur through actual extension maps.
The map fixing $v,w$ and sending $X$ to $X^{-1}$
is an order-two automorphism of $E_0$.
On the kernel it is precisely $\zeta_0$ in
the statement.
It commutes with $\operatorname{Inn}(v)$ and
\[
                 \zeta_0(c_0)=X^{-1}c_0^{-1}X,
\]
so it is peripheral of multiplier $-1$.
The two commuting involutions have respective
homology matrices
$\operatorname{diag}(1,-1)$ and
$\operatorname{diag}(-1,1)$.
They give four distinct outer classes.
This proves (D.64b).
A normalizer of an order-two subgroup fixes its
sole nonidentity element and is its centralizer.
Finally equality of the two orbit images of
$\theta,\theta'$ is exactly the statement that
$\theta^{-1}\theta'$ belongs to this computed
centralizer, proving the stated actual orbit fibres.

\emph{The exact multiplier-one conclusion in the one-class model.}
The automorphisms $\alpha_1$ and $\zeta_1$ commute
on both free generators and have order two.
Moreover
\[
        \zeta_1(c_0)=(XY)^{-1}c_0(XY),
        \qquad (\zeta_1)_*=-I .
                                                               \tag{D.64l}
\]
Thus $\zeta_1$ has multiplier one and fixed-line
sign $-1$, while $b_1$ has multiplier $-1$
and fixed-line sign $1$.
They provide the four indicated elements.
The injective characters in (D.64c) allow at
most two elements with multiplier one.
Multiplication by $b_1$ gives exactly the two
with multiplier $-1$.
This proves (D.64d).
Every torsion element of this abelian centralizer
generates a finite subgroup; together with the
displayed four-element subgroup it still generates
a finite group.
D.63g bounds that group by order four.
Consequently the displayed four elements are
exactly all its torsion.
This argument establishes the stated character
embedding and its torsion subgroup. D.65 below
computes its remaining multiplier image through
the specified cone-filling quotients.

The program
\path{code/part_c/standard_negative_centralizers_d64.py}
checks the exact free-group identities in
(D.64a), (D.64g), and (D.64l), and checks both
dihedral transporter formulas over a finite range.
The global conclusion (D.64k) is the proved
all-modulus statement above.
\end{proof}
